\chapter{ต้นไม้ตัดสินใจและการสุ่มป่าไม้}
ต้นไม้ตัดสินใจ (Decision Tree) เป็นอัลกอริทึมการเรียนรู้ของเครื่องที่มีความสำคัญและได้รับความนิยมอย่างกว้างขวางในการแก้ปัญหาทั้งการจำแนกประเภท (Classification) และการทำนายค่า (Regression) เนื่องจากมีโครงสร้างที่เข้าใจง่าย สามารถตีความผลได้ชัดเจน และไม่ต้องการการปรับแต่งข้อมูลมากนัก อย่างไรก็ตาม ต้นไม้ตัดสินใจเดี่ยวมักมีปัญหา Overfitting และมีความไวต่อการเปลี่ยนแปลงของข้อมูล ซึ่งนำไปสู่การพัฒนาเทคนิค Ensemble Learning โดยเฉพาะการสุ่มป่าไม้ (Random Forest) ที่เป็นการรวมต้นไม้ตัดสินใจหลายต้นเข้าด้วยกันผ่านกระบวนการ Bagging และการสุ่มเลือกตัวแปรย่อย (Feature Sampling) เพื่อเพิ่มความแม่นยำและลดปัญหา Overfitting ทำให้ Random Forest กลายเป็นหนึ่งในอัลกอริทึมที่ทรงพลังและใช้งานได้จริงในปัญหาหลากหลายสาขา ตั้งแต่การวิเคราะห์ข้อมูลทางการแพทย์ การเงิน ไปจนถึงการจำแนกภาพและการประมวลผลภาษาธรรมชาติ ในบทนี้จะกล่าวถึงโครงสร้างและหลักการทำงานของต้นไม้ตัดสินใจ กระบวนการแบ่งข้อมูลโดยใช้แนวคิด Entropy และ Information Gain Ratio การตัดแต่งกิ่ง (Pruning) เพื่อป้องกัน Overfitting รวมถึงหลักการของ Random Forest และเทคนิค Ensemble Learning ที่นำไปสู่ประสิทธิภาพที่ดีขึ้น


\section{ต้นไม้ตัดสินใจ (Decision Trees)}
\subsection{โครงสร้างและหลักการทำงาน}
\subsection{เอนโทรปี}
\subsection{เกรนความรู้และอัตราเกรนความรู้}


\section{การสุ่มป่าไม้ (Random Forest)}
\subsection{แนวคิด Ensemble Learning}
\subsection{Bagging และ Bootstrap Sampling}
\subsection{ความสำคัญของตัวแปร (Feature Importance)}



\newpage
%\RefChapter